\section{Proof of Theorem~\ref{thm:main}}\label{app:proof}
Write $\gamma=\eta\tau$ and let $\hat g_t$ be the aggregated direction, so that $\mathbf w^{t+1}=\mathbf w^t-\gamma\hat g_t$. All expectations are conditional on $\mathbf w^t$ unless stated otherwise; the participation and depth process is independent of the iterate by assumption.

\paragraph{Step 1: bias decomposition.}
For client $k$ with depth $d$, let $g_k=\frac1\tau\sum_{i<\tau}\tilde\nabla F_k^{(d)}(\mathbf w_{k,i})$ be the average of its local stochastic gradients along the local trajectory $\mathbf w_{k,0}=\mathbf w^t$. By (A2) $\|\mathbf w_{k,i}-\mathbf w^t\|\le\eta iG$. Since stochastic gradients are unbiased given the local iterate and $F_k^{(d)}$ is $L_s$-smooth,
\[
\big\|\mathbb E[g_k]-\nabla F_k^{(d)}(\mathbf w^t)\big\|\le\frac1\tau\sum_{i<\tau}L_s\eta iG=\frac{L_s\eta G(\tau-1)}2\le\frac{L_sG\gamma}2=:\varepsilon .
\]
Quantisation is unbiased (A3), so it does not change the mean. For block $l$ the aggregate with weights $p_k/\pi_{k,l}$ has mean
\[
\mathbb E[\hat g_{t,l}]=\sum_kp_k\,\mathbb E_d\big[\nabla_lF_k^{(d)}(\mathbf w^t)\mid d\ge l,k\in\mathcal S\big]+\mathrm{drift}_l ,
\]
and by Minkowski's inequality the stacked drift over all blocks has norm at most $\varepsilon$.

\paragraph{Step 2: truncation bias.}
Since $\nabla_lF_k^{(d)}=\sum_{j=l}^{d}\lambda_j\nabla_lf_k^{j}$,
\[
\mathbb E_d\big[\nabla_lF_k^{(d)}\mid d\ge l\big]=\sum_{j\ge l}\lambda_j\Pr[d\ge j\mid d\ge l]\,\nabla_lf_k^{j}.
\]
For $j>l$ we have $\{d\ge j\}\subseteq\{d\ge l\}$, hence $\Pr[d\ge j\mid d\ge l]=\sigma_{k,j}/\sigma_{k,l}\ge\sigma_{k,j}\ge\sigma_j$, while the $j=l$ term has probability one. Comparing with $\nabla_lF_k=\sum_{j\ge l}\lambda_j\nabla_lf_k^{j}$,
\[
\Big\|\nabla_lF_k-\mathbb E_d[\nabla_lF_k^{(d)}\mid d\ge l]\Big\|\le G_l\sum_{j>l}\lambda_j(1-\sigma_j),
\]
and averaging over $k$ with weights $p_k$ gives block bias at most $B_l$. Stacking the blocks, $\mathbb E[\hat g_t]=\nabla F(\mathbf w^t)-b_t+e_t$ with $\|b_t\|\le B$ and $\|e_t\|\le\varepsilon$.

\paragraph{Step 3: descent.}
By $L_s$-smoothness of $F$, $F(\mathbf w^{t+1})\le F(\mathbf w^t)-\gamma\langle\nabla F(\mathbf w^t),\hat g_t\rangle+\frac{L_s\gamma^2}{2}\|\hat g_t\|^2$. Using $\langle a,a-c\rangle\ge\frac12\|a\|^2-\frac12\|c\|^2$,
\[
\mathbb E\langle\nabla F,\hat g_t\rangle\ge\tfrac12\|\nabla F\|^2-\tfrac12(B+\varepsilon)^2 ,
\]
and by (A4) $\mathbb E\|\hat g_t\|^2\le(\|\nabla F\|+B+\varepsilon)^2+\sigma_h^2\le2\|\nabla F\|^2+2(B+\varepsilon)^2+\sigma_h^2$. Therefore
\[
\mathbb EF(\mathbf w^{t+1})\le F(\mathbf w^t)-\gamma\big(\tfrac12-L_s\gamma\big)\|\nabla F\|^2+\gamma\big(\tfrac12+L_s\gamma\big)(B+\varepsilon)^2+\tfrac{L_s\gamma^2}{2}\sigma_h^2 .
\]
For $\gamma\le1/(4L_s)$ the coefficients are at least $\frac14$ and at most $\frac34$. Taking total expectation, summing over $t<T$ and dividing by $\gamma T/4$,
\[
\frac1T\sum_t\mathbb E\|\nabla F(\mathbf w^t)\|^2\le\frac{4\Delta}{\gamma T}+3(B+\varepsilon)^2+2L_s\gamma\sigma_h^2 .
\]

\paragraph{Step 4: constants.}
Using $(B+\varepsilon)^2\le2B^2+2\varepsilon^2$ and $\varepsilon\le L_sG\gamma/2$, we get $3(B+\varepsilon)^2\le6B^2+\frac32L_s^2G^2\gamma^2\le6B^2+\frac38L_sG^2\gamma\le6B^2+2L_s\gamma\frac{G^2}4$, where the second inequality uses $\gamma\le1/(4L_s)$. Hence the bound is at most $\frac{4\Delta}{\gamma T}+2L_s\gamma\tilde\sigma^2+6B^2$ with $\tilde\sigma^2=\sigma_h^2+G^2/4$. If $\gamma=\sqrt{\Delta/(L_s\tilde\sigma^2T)}\le1/(4L_s)$ the first two terms sum to $6\sqrt{L_s\tilde\sigma^2\Delta/T}$. Otherwise $\gamma=1/(4L_s)$ and $\sqrt{\Delta/(L_s\tilde\sigma^2T)}>1/(4L_s)$, so the first term equals $16L_s\Delta/T$ and the second is at most $2\sqrt{L_s\tilde\sigma^2\Delta/T}$. In both cases \eqref{eq:thm} holds. \qed

\paragraph{Corollary~\ref{cor:cov}.}
With $\lambda_j=1/L$, $G_l=G$ and $\sigma_j=\sigma$, $B_l=G(1-\sigma)(L-l)/L$ and $B^2=G^2(1-\sigma)^2\sum_{l=1}^{L}\big(\tfrac{L-l}{L}\big)^2=G^2(1-\sigma)^2\tfrac{(L-1)(2L-1)}{6L}$, so $6B^2=\tfrac{(L-1)(2L-1)}{L}G^2(1-\sigma)^2$. Requiring this floor to be at most $\epsilon/2$ gives \eqref{eq:cov}.
